% Notas de aula — Capítulo 15: Perigos de Tempestades
% Curso: Meteorologia Prática (baseado em R. Stull, Practical Meteorology, CC BY-NC-SA 4.0)
% Caixas verdes = conteúdo literal do quadro; linhas verdes = encenação.
\documentclass[11pt]{article}
\usepackage{../comum/preambulo}
\graphicspath{{../figuras/cap15/}}

\title{Capítulo 15 --- Perigos de Tempestades\\
{\large Notas de aula (roteiro de quadro-negro)}}
\author{Curso de Meteorologia Prática}
\date{}

\begin{document}
\maketitle
\thispagestyle{fancy}

\noindent\textbf{Plano sugerido:} 2 aulas de 2\,h.
\emph{Aula 1:} \S\ref{sec:granizo}--\S\ref{sec:down} (granizo,
downbursts e rajadas).
\emph{Aula 2:} \S\ref{sec:raios}--\S\ref{sec:tornado} (relâmpagos,
tornados e a escala EF).

\section{A máquina e seus produtos}

\begin{noquadro}[abertura]
\textbf{Cap.~15 --- Perigos de Tempestades}\\[2pt]
a supercélula do Cap.~14 é a fábrica; os produtos:\\
granizo \quad $|$ \quad downburst/rajadas \quad $|$ \quad relâmpago
\quad $|$ \quad tornado\\[2pt]
cada perigo $=$ um pedaço de física clássica levado ao extremo.
\end{noquadro}

Para físicos, este capítulo é um parque de diversões: balanço
peso--arrasto (granizo), termodinâmica de evaporação (downburst),
eletrostática (raios) e equilíbrio ciclostrófico (tornado)
\javimos{10}{ciclostrófico}.

\section{Granizo}\label{sec:granizo}

\quadro{Perguntar: o que segura uma pedra de 5\,cm a 8\,km de altura?
Montar o balanço peso $=$ arrasto no quadro.}

\begin{formadif}[o updraft mínimo (T1; Caso~1)]
Peso $=$ arrasto:
\[ \rho_{gelo}\,g\,\frac{\pi D^3}{6} =
C_d\,\rho_{ar}\,\frac{\pi D^2}{8}\,w_T^2
\quad\Longrightarrow\quad
w_T = \sqrt{\frac{4\,g\,\rho_{gelo}\,D}{3\,C_d\,\rho_{ar}}}
\propto \sqrt{D}. \]
Pedra de 5\,cm $\Rightarrow$ $w \gtrsim 32$\,m/s: granizo grande é
\emph{diagnóstico} de supercélula \javimos{14}{updrafts}.
\end{formadif}

\begin{noquadro}[a biografia da pedra]
embrião (graupel, Cap.~7) $\to$ voos repetidos pelo updraft $\to$
camadas alternadas (crescimento seco/úmido) $\to$ cai quando
$D$ vence o updraft\\[2pt]
derretimento na queda: nível de congelamento alto derrete o granizo
pequeno (N1) — por isso o equador tem menos granizo que os
subtrópicos.
\end{noquadro}

O crescimento por acreção é o \emph{riming} levado ao extremo
\javimos{7}{riming}; as camadas contam a história térmica dos voos. No
radar: 55$+$\,dBZ e assinaturas de polarização dual
\javimos{8}{dual-pol}.

\section{Downbursts e rajadas}\label{sec:down}

\quadro{Desenhar a coluna de chuva evaporando, o ar frio desabando e o
jato se espalhando no chão. Comparar com o jato de pia invertido.}

\begin{formadif}[DCAPE (T2; Caso~2)]
O espelho descendente do CAPE:
\[ \mathrm{DCAPE} = \int_0^{z_0} g\,
\frac{T_{v,e} - T_{v,p}}{T_{v,e}}\,dz,
\qquad w_{down} = \sqrt{2\,\mathrm{DCAPE}}. \]
A parcela desce pelo seu $\theta_w$ \javimos{4}{bulbo úmido},
resfriada pela evaporação da chuva; DCAPE de 500--1500\,J/kg dá
rajadas de 30--55\,m/s.
\end{formadif}

\begin{noquadro}[seco × úmido]
camada sub-nuvem SECA $\Rightarrow$ evaporação vigorosa $\Rightarrow$
downburst forte (N2)\\
downburst \textbf{seco}: cai de virga, quase sem eco no radar — a
armadilha da aviação\\
no chão: microburst ($<$4\,km) cisalha a aproximação final — os
acidentes clássicos dos anos 70--80
\end{noquadro}

A rajada espalhada é a frente de rajada \javimos{14}{piscina fria} —
que planta a próxima célula. Escala maior: \emph{derechos}, corredores
de rajadas de linhas de instabilidade longevas \javimos{14}{RKW}.

\section{Relâmpagos}\label{sec:raios}

\quadro{Montar o circuito no quadro: gerador (colisões
graupel--cristal no updraft), capacitor (nuvem--solo), resistor de
fuga, e a ruptura dielétrica como chave.}

\begin{formadif}[o capacitor da tempestade (Caso~3)]
\[ \frac{dQ}{dt} = I_c - \frac{Q}{RC},\qquad
E = \frac{Q}{C\,d} \to E_{rup} \simeq 3\times10^5\ \mathrm{V/m}
\ (\text{com gotas}). \]
Dente de serra: carrega em dezenas de segundos, descarrega
$\sim$85\,\% num raio de $\sim$5\,C e $\sim$1--10\,GJ. A taxa de raios
cresce brutalmente com o updraft (N3) — o mapa global de relâmpagos é
um mapa de convecção continental profunda.
\end{formadif}

\begin{noquadro}[números de segurança]
flash-to-bang: som a 343\,m/s $\Rightarrow$ \textbf{3\,s por km}
(T3)\\
regra 30--30: trovão em $<$30\,s do clarão $\Rightarrow$ abrigo; sair
só 30\,min após o último trovão\\
canal a $\sim$30\,000\,K (5$\times$ a superfície do Sol) — o trovão é
o estrondo da expansão
\end{noquadro}

A separação de carga exige gelo $+$ água super-resfriada $+$ updraft
— exatamente a zona de crescimento do granizo \javimos{7}{fase mista}:
raio e granizo são irmãos.

\section{Tornados}\label{sec:tornado}

\quadro{Desenhar o vórtice de Rankine ($v \propto r$ dentro,
$\propto 1/r$ fora) e integrar a queda de pressão no quadro.}

\begin{formadif}[Rankine ciclostrófico (T4; Caso~4)]
Sem Coriolis nessa escala \javimos{10}{quem sente Coriolis}:
\[ \frac{dP}{dr} = \frac{\rho v^2}{r}
\quad\Longrightarrow\quad
\Delta P_{centro} = \rho\,v_{max}^2
\ \ (\text{metade do núcleo, metade de fora}). \]
EF3 (70\,m/s): $\sim$50\,hPa em 100\,m; EF5 (110\,m/s):
$\sim$120\,hPa — gradientes que nenhuma escala sinótica alcança.
\end{formadif}

\begin{noquadro}[da rotação da tempestade ao chão]
mesociclone (Cap.~14, hodógrafo curvo) $\to$ estiramento sob o updraft
($-(f{+}\zeta)D$! \javimos{13}{esticamento}) $\to$ funil\\
escala EF: dano $\to$ vento; carga $q = \tfrac12\rho v^2$ (T5) —
dobrar $v$ quadruplica a força\\
HS: giro horário; com translação, o lado \textbf{esquerdo} da
trajetória é o pior (N4)
\end{noquadro}

Tornados não-supercelulares (landspouts, trombas d'água) nascem de
cisalhamento raso esticado por cumulus congestus — mais fracos, mas
frequentes no litoral brasileiro.

\section{Síntese da aula}

\begin{noquadro}[mapa final --- canto do quadro]
\[ w_T = \sqrt{\tfrac{4g\rho_i D}{3C_d\rho_a}} \quad
w_{down} = \sqrt{2\,\mathrm{DCAPE}} \quad
\tfrac{dQ}{dt} = I_c - \tfrac{Q}{RC} \quad
\Delta P = \rho v_{max}^2 \]
quatro perigos, quatro pedaços de física clássica no extremo
\end{noquadro}

\section{Exercícios complementares}\label{sec:exercicios}

Soluções (professor) em \texttt{cap15\_solucoes.ipynb}.

\subsection*{Teóricos}

\begin{enumerate}[label=\textbf{T\arabic*.},itemsep=4pt]
  \item Do balanço peso--arrasto, deduza
        $w_T = \sqrt{4g\rho_{gelo}D/(3C_d\rho_{ar})}$ e calcule o
        updraft mínimo para pedras de 1, 2{,}5, 5 e 10\,cm.
  \item Defina o DCAPE como espelho do CAPE, explicando por que a
        parcela de downdraft desce pelo seu $\theta_w$ e qual o papel
        da umidade da camada sub-nuvem.
  \item Justifique a regra dos ``3 segundos por km'' do flash-to-bang
        e estime a energia de um raio a partir de $Q \sim 5$\,C e
        $V \sim 10^8$\,V.
  \item Integre o equilíbrio ciclostrófico para o vórtice de Rankine e
        mostre que $\Delta P = \rho v_{max}^2$, com metade da queda
        dentro do núcleo.
  \item Calcule a pressão dinâmica $q = \tfrac12\rho v^2$ para os
        limiares EF0--EF5 e discuta a não-linearidade da destruição.
\end{enumerate}

\subsection*{Numéricos (no notebook)}

\begin{enumerate}[label=\textbf{N\arabic*.},itemsep=4pt]
  \item Integre a queda com derretimento e ache o diâmetro mínimo que
        sobrevive para níveis de congelamento a 3, 4 e 5\,km.
  \item Faça o déficit evaporativo depender da UR sub-nuvem e trace a
        rajada final; por que downbursts secos são traiçoeiros?
  \item Varra a corrente de carga do capacitor e trace a taxa de
        raios; conecte com a sensibilidade $\sim w^6$ observada.
  \item Some translação ao Rankine e mapeie o vento total; qual lado
        da trajetória concentra o dano no HS?
\end{enumerate}

\vspace{1em}
\noindent\rule{\linewidth}{0.4pt}\\
{\scriptsize Material derivado de R.~Stull, \emph{Practical Meteorology}
(CC BY-NC-SA 4.0); este documento herda a mesma licença.}

\end{document}
